%%%PAGE 171 rot=0 legibility=good summary=力的平行四边形实验（①两测力计+一测力计 ②橡皮筋 ③三力平衡转+相似+做圆，A-D选项含红笔批改）；弹簧k误差一句；测污水电阻率（大R内接法、分压）
%%%BLOCK id=171-01 ch=02 sec=验证力的平行四边形定则·实验 kind=method exp=1 alt=none cont=none y=0.02-0.22
\textbf{力的平行四边形}

①

%FIG p171_a 0.14 0.043 0.47 0.185
\notefig[0.4]{figures/p171_a.png}

记录$F$大小与方向

$F_3$合力：橡皮/测力计
%%%END
%%%BLOCK id=171-02 ch=02 sec=验证力的平行四边形定则·实验 kind=method exp=1 alt=none cont=none y=0.23-0.43
② 橡皮筋

%FIG p171_b 0.235 0.258 0.49 0.402
\notefig[0.35]{figures/p171_b.png}

在力的方向上画伸长量

%FIG p171_c 0.49 0.35 0.63 0.41
\notefig[0.2]{figures/p171_c.png}

\uwave{不用让$O$在同一位置了}
%%%END
%%%BLOCK id=171-03 ch=02 sec=验证力的平行四边形定则·实验 kind=method exp=1 alt=none cont=none y=0.43-0.68
③ 三力 $\to$ 转$+$相似$+$做圆

%FIG p171_d 0.14 0.44 0.51 0.64
\notefig[0.4]{figures/p171_d.png}

% 红笔批改：A 行末有红勾；B 行末红×并补写「和方向」；C 行末有红勾/红线并旁注「用细绳标方向」；D「大一些」下有红线+红勾
A：$O$点不能动；同次\ \red{$\checkmark$}

B：记下测力计读数，$O$点位置\red{$\times$和方向} % 「记下」为行间小字插入

C：拉橡皮筋的细绳长一些，效果好\ \red{$\checkmark$}\quad \red{用细绳标方向}

D：拉力要适当大一些\ \red{$\checkmark$}
%%%END
%%%BLOCK id=171-04 ch=02 sec=弹簧·胡克定律实验 kind=note exp=1 alt=none cont=none y=0.68-0.73
结点$D$偏高没有影响，直$R$斜\unclear{放}$k$偏小 % CHECK: 「D」疑为「O」；「直R」疑为「直尺」；该句疑指弹簧(k)实验中直尺斜放读数偏大 → k偏小
%%%END
%%%BLOCK id=171-05 ch=10 sec=测电阻·内接法与分压 kind=method exp=1 alt=none cont=none y=0.75-0.93
测污水电阻率\quad 大$R$\ 内接法\quad 选小滑变

%FIG p171_e 0.115 0.777 0.245 0.885
\notefig[0.25]{figures/p171_e.png}

\[
\Bigl(I-\frac{U}{R_V}\Bigr)R_x=U
\]

\[
U=\frac{R_xR_V}{R_x+R_V}\,I
\]

% 页面右侧另有两处潦草涂改的小式，疑为分式「U/R_x +」与「U/R_V +」
\unclear{$\dfrac{U}{R_x}+$}\quad\unclear{$\dfrac{U}{R_V}+$}
%%%END
%%%PAGEEND 171
